\section{问题重述}

\subsection{问题背景}

中药材的干燥是决定成品品质的关键工序之一，热风烘干因设备简单、易于控制而成为常用方式。
整个过程分为预热平衡与恒温干燥两个阶段：预热阶段药材由室温逐步升温、内部水分开始向表面迁移；
恒温干燥阶段烘房维持稳定的温湿环境，水分持续从表面蒸发、内部水分不断向外补充，直至达到干燥要求。
工艺参数选取不当会带来干燥效率低、能耗高、成品品质不稳定等问题，而传统的逐批试验优化成本高、周期长，
因此借助机理分析与数值仿真获得药材内部温度与水分浓度的时空规律，具有重要的工程意义。

在实际干燥过程中，药材会因水分流失而收缩，水分扩散路径随之缩短；与此同时，扩散系数本身又受温度和浓度影响。二者对干燥时长的影响方向相反、相互竞争。因此，只有通过建模进行定量分析，才能判断哪一效应占主导。

\subsection{问题提出}

题目围绕圆柱形中药材在热风烘干过程中的温度场与水分浓度场提出四个问题：

\begin{itemize}[label=\smash{\large\ding{226}}, leftmargin=2.4em, itemsep=3pt]
  \item \textbf{问题一}：给出烘干初期 $0\sim1800$~s 内药材内部温度与水分浓度的变化规律，
        并按题目规定的时刻与位置列表。同时输出该时段内每隔固定时间、固定径向间隔的完整计算结果。
  \item \textbf{问题二}：在恒温环境下建立整个烘干过程中温湿场变化的模型，
        按题目规定的时刻与位置给出结果，并输出完整计算结果。
  \item \textbf{问题三}：依据题目给定的干燥要求确定恒温干燥所需的时间，
        并按题目规定的时刻与位置给出水分浓度。
  \item \textbf{问题四}：在药材随失水而收缩的条件下确定烘干时长，
        并按题目规定的时刻与位置给出水分浓度，输出完整计算结果。
\end{itemize}
